\section{Smooth Manifolds}
\begin{definition}\label{4.1}
    Let $X\subseteq \mathbb{R}^k,Y\subseteq \mathbb{R}^l$ be arbitrary subsets. A map $f:X\to Y$ is called \textbf{smooth}\index{smooth}, if for each $x\in X$ 
    there exists an open neighborhood $U\subseteq \mathbb{R}^k$ of $x$ and a smooth map $F:U\to \mathbb{R}^l$ such that $F|_{U\cap X}=f$.\\
    $f$ is called a \textbf{diffeomorphism}\index{diffeomorphism}, if $f$ is a homeomorphism and both $f$ and $f^{-1}$ are smooth.\\  
\end{definition}
\begin{proposition}\label{4.2}
    $f:X\to Y;g:Y\to Z$ smooth $\implies g\circ f:X\to Z$ smooth.\\
\end{proposition}
\begin{definition}\label{4.3}
    $M\subseteq \mathbb{R}^k$ is called a \textbf{smooth $n$-dimensional manifold}\index{manifold}, if each $x\in M$ has an open neighborhood 
    $W\subseteq \mathbb{R}^k$ such that $W\cap M$ is diffeomorphic to an open subset $U\subseteq \mathbb{R}^n$.\\
    Any particular diffeomorphism $g:U\to W\cap M$ is called a \textbf{parametrization}\index{parametrization} of $W\cap M$. \\
    By definition $M$ is a manifold of dimension $0$ if each $x\in M$ has an open neighborhood $W\subseteq \mathbb{R}^k$ such that $W\cap M=\{x\}$.  
\end{definition}
\begin{example}\label{4.4}
    $\mathbb{S}^2=\{(x,y,z)|x^2+y^2+z^2=1\}$ is a smooth $2$-dimensional manifold.\\
    We see that $(x,y)\to(x,y,\sqrt{1-x^2-y^2})$ for $x^2+y^2<1$ is a parametrization of the upper hemisphere ($\mathbb{S}^2\cap\{z>0\}$).\\
    We can define similar parametrizations for $\{x>0\},\{y<0\},\{y>0\},\{z<0\}$ and $\{x<0\}$.\\ 
\end{example}
\begin{remark}\label{4.5}
    Let $U\subseteq \mathbb{R}^k$ and $f:U\to \mathbb{R}^l$ be smooth.\\
    $df_x:\mathbb{R}^k\to \mathbb{R}^l$ is a linear map, defined by $df_x(h)=\lim_{t\to 0}\frac{f(x+th)-f(x)}{t}$.\\ 
    We remember:
    \begin{itemize}
        \item chain rule: $d(g\circ f)_x=dg_{f(x)}\cdot df_x$.
        \item $dId_U=Id_{\mathbb{R}^k}$.
        \item $L:\mathbb{R}^k\to \mathbb{R}^l$ linear $\implies dL_x=L$.
    \end{itemize}
    In particular if $f:U\subseteq \mathbb{R}^k\to V\subseteq \mathbb{R}^l$ is a diffeomorphism, then $k=l$ and $df_x$ is non-singular.\\
    The other way around is also true: If $df_x$ is non-singular for all $x\in U$, then via the inverse function theorem $f$ maps any sufficiently small open set $U'$ around $x$ diffeomorüphically onto an open set $f(U')$.\\  
    Now let $M\subseteq \mathbb{R}^k$ be a smooth $m$-dimensional manifold and $x\in M$.\\
    Let $g:U\subseteq \mathbb{R}^m\to M\subseteq\mathbb{R}^k$ be a parametrization, and let $g(U)$ be a neighborhood of $x\in M$.\\
    Then $dg_u:\mathbb{R}^m\to \mathbb{R}^k$ is well-defined $\forall u\in U$.\\
    We define the \textbf{tangent space}\index{tangent space} of $M$ at $x$ as $T_xM:=dg_u(\mathbb{R}^m)$ such that $g(u)=x$.\\
    $T_xM$ is well-defined, because if $h:V\subseteq \mathbb{R}^m\to M$ is another parametrization with $g(V)$ being a neighborhood of $x\in M$. Take some $v\in V$ such that $h(v)=x$.\\
    Then $h^{-1}\circ g$ maps a neighbourhood $U_1$ of $u$ diffeomorphically onto a neighborhood $V_1$ of $v$.\\   
    \begin{center}
    \begin{minipage}{0.48\textwidth}
    \centering
    \begin{tikzcd}[column sep=small, row sep=large]
    & \mathbb{R}^l & \\
    U_1 \arrow[rr, "h^{-1}\circ g"'] \arrow[ur, "g"] & & V_1 \arrow[ul, "h"']
    \end{tikzcd}
    \end{minipage}
    \hfill
    \begin{minipage}{0.48\textwidth}
    \centering
    \begin{tikzcd}[column sep=small, row sep=large]
    & \mathbb{R}^k & \\
    \mathbb{R}^m \arrow[rr, "d(h^{-1}\circ g)_u"'] \arrow[ur, "dg_u"] & & \mathbb{R}^m \arrow[ul, "dh_v"']
    \end{tikzcd}    
    \end{minipage}
    \end{center}
    This implies that $Im(dg_u)=Im(dh_v)$.\\
    $T_xM$ is an $m$-dimensional vector space ($rank(dg_u)=m$):
    $g^{-1}:g(U)\to U$ is smooth, hence we may choose $x\in W$ open in $\mathbb{R}^k$ such that $F:W\to \mathbb{R}^m$ is smooth and $F|_{W\cap g(U)}=g^{-1}$.\\
    \begin{center}
\begin{minipage}{0.48\textwidth}
\centering
\begin{tikzcd}[column sep=small, row sep=large]
& W \arrow[dr, "F"] & \\
U \arrow[ur, "g"] \arrow[rr, "inclusion"'] & & \mathbb{R}^m
\end{tikzcd}
\end{minipage}
\hfill
\begin{minipage}{0.48\textwidth}
\centering
\begin{tikzcd}[column sep=small, row sep=large]
& \mathbb{R}^k \arrow[dr, "dF_x"] & \\
\mathbb{R}^m \arrow[ur, "dg_u"] \arrow[rr, "\mathrm{id}"'] & & \mathbb{R}^m
\end{tikzcd}
\end{minipage}
    \end{center}
\end{remark}
\begin{example}\label{4.6}
    \begin{itemize}
        \item Let $M=\{(x,y,z)\in\mathbb{R}^3|z=x^2+y^2\}$. Then
        \[g:\mathbb{R}^2\to \mathbb{R}^3,\quad g(x,y)=(x,y,x^2+y^2)\]
        is a parametrization of $M$.\\
        We calculate $T_{(x,y,z)}M=IM(dg_{(x,y)})=span\{(1,0,2x),(0,1,2y)\}$.\\
        \item Let $M=\mathbb{S}^2\subseteq \mathbb{R}^3$. Then
        \[g(x,y)=(x,y,\sqrt{1-x^2-y^2})\]
        is a parametrization of the upper hemisphere.\\
        \item Let $M=\mathbb{T}^2\subseteq \mathbb{R}^3$ be the torus. We can parametrize $M$ by first parametrizing a circle in the $x-z$-plane ($c(t)=(R+r\cos t,0,r\sin t)$) and then rotating this circle around the $z$-axis.\\
        This gives us the parametrization
        \[g(s,t)=((R+r\cos t)\cos s,(R+r\cos t)\sin s,r\sin t)\]
        We cover $M$ by using $3$ charts: $(0,2\pi)^2,(\frac{\pi}{2},\frac{5}{2}\pi)^2,(-\pi,\pi)^2$.\\
        We calculate $T_{g(s,t)}M=Im(dg_{(s,t)})=span\{\begin{pmatrix}
            -(R+r\cos t)\sin s\\(R+r\cos t)\cos s\\
            0
        \end{pmatrix},\begin{pmatrix}
            -r\sin t\cos s\\-r\sin t\sin s\\r\cos t
        \end{pmatrix}\}$.
    \end{itemize}
\end{example}
\begin{definition}\label{4.7}
    Now consider $f:M\to N$ smooth, where $M\subseteq \mathbb{R}^k$ and $N\subseteq \mathbb{R}^l$ are smooth manifolds.\\
    We define $df_x:T_xM\to T_{f(x)}N$ as follows:\\
    Take $F:W\to \mathbb{R}^l$ smooth such that $F|_{W\cap M}=f$ where $W$ is an open neighborhood of $x$ in $\mathbb{R}^k$.\\
    We set $df_x(v):=dF_x(v)$.\\
    Now we need to show that $df_x(v)\in T_{f(x)}N$ and that $df_x$ is independent of the choice of $F$.\\  
    Let us choose a parametrization $g:U\to M\cap W\subseteq \mathbb{R}^k$ and $h:V\to N\subseteq\mathbb{R}^l$.\\
    Then $dF_x(T_xM)=dF_x(dg_u)\subseteq dh_v=T_{f(x)}N$.\\
    Also $dF_x=dh_v\circ d(h^{-1}\circ f\circ g)_u\circ dg_u^{-1}$, hence $df_x$ is independent of the choice of $F$.\\
    \begin{itemize}
        \item $df_x$ is linear.
        \item the chain rule holds.
        \item $dId_M=Id_{T_xM}$.
    \end{itemize}
    In particular if $f:M\to N$ is a diffeomorphism, then $df_x:T_xM\to T_{f(x)}N$ is an isomorphism so $dim M=dim N$.\\
    \begin{center}
\begin{minipage}{0.48\textwidth}
\centering
\begin{tikzcd}[column sep=large, row sep=large]
U \arrow[r, "g"] \arrow[d, "h^{-1}\circ f\circ g"'] & W \arrow[d, "F"] \\
V \arrow[r, "h"'] & \mathbb{R}^l
\end{tikzcd}
\end{minipage}
\hfill
\begin{minipage}{0.48\textwidth}
\centering
\begin{tikzcd}[column sep=large, row sep=large]
\mathbb{R}^m \arrow[r, "dg_u"] \arrow[d, "d(h^{-1}\circ f\circ g)_u"'] & \mathbb{R}^k \arrow[d, "dF_x"] \\
\mathbb{R}^m \arrow[r, "dh_v"'] & \mathbb{R}^l
\end{tikzcd}
\end{minipage}
\end{center}
\end{definition}